% !TeX root = ../../Book3.tex
% 纲要标题：### 7.4 曲线分支与非约化结构
% 纲要位置：工作记录/计划纲要.md，第 2127 行。
% 纲要细目（写作范围）：
% * 尖点、节点、交叉直线与加厚点；经典约化模型与非约化测试代数的区别
% * 约化、不可约与正规性的区别；局部整闭性及素谱与经典参数映射的同胚
% * 不用图像替代局部环计算；导子零点集之外的局部同构、余切空间与一阶切空间
% ---

\section{曲线分支与非约化结构}
\label{b3:sec:0704}

尖点、节点与交叉直线都属于经典仿射代数集，具有约化坐标环；加厚点 \(k[\varepsilon]/(\varepsilon^2)\) 则是非约化有限型代数，不能作为普通代数集的坐标环。它用来检验域值点会遗漏什么信息。固定代数闭域 \(k\)，节点例额外要求 \(\operatorname{char}k\ne2\)。本节使用第六章已经计算的正规化与导子，比较这四个模型的纤维代数和局部环。

尖点与节点的计算使用\cref{b3:prop:cusp-normalization,b3:prop:node-normalization}的显式环包含及\secref{b3:sec:0606}的导子计算；交叉直线的全分式环与正规化已在\secref{b3:sec:0604}通过常数项相等的多项式对具体求出。

\subsection{尖点、节点、交叉直线与加厚点}
\label{b3:subsec:0704:01}

尖点 \(C=Z_k(y^2-x^3)\) 的坐标环由\cref{b3:prop:cusp-normalization}识别为 \(A=k[t^2,t^3]\)，正规化为 \(B=k[t]\)。参数映射 \(t\mapsto(t^2,t^3)\) 双射：原点来自 \(t=0\)，其余点 \((a,b)\) 满足 \(a\ne0\)，由 \(t=b/a\) 唯一恢复参数，其平方、立方恰为 \(a,b\)。环包含却严格。\cref{b3:prop:cusp-conductor}给出导子和商模；原点的极大理想 \(\mathfrak m=(t^2,t^3)_A\) 在 \(B\) 中生成 \((t^2)\)，因为 \(t^3=t\cdot t^2\)。于是
\[
\mathfrak c=t^2k[t],\qquad B/A\cong k\cdot[t],
\qquad B\otimes_Ak=B/\mathfrak mB\cong k[t]/(t^2).
\]
其中 \([t]=t+A\) 张成一维商空间 \(B/A\)。所以即使点的个数没有改变，原点的纤维代数仍有非零幂零元；集合纤维只有一个点，其约化坐标环为 \(k\)。原坐标环也确实少了一个函数 \(t\)。

\ProseParagraphBreak
节点 \(N=Z_k\bigl(y^2-x^2(x+1)\bigr)\) 使用参数 \(x=t^2-1,y=t(t^2-1)\)。由\cref{b3:prop:node-normalization}和\cref{b3:prop:node-conductor}，坐标环与正规化为
\[
A=k+(t^2-1)k[t]\subseteq B=k[t],\qquad
\mathfrak c=(t^2-1)k[t].
\]
原点理想 \(\mathfrak m=(x,y)_A\) 在 \(B\) 中生成 \((t^2-1)\)，所以纤维代数为 \(B/\mathfrak mB=k[t]/(t^2-1)\cong k\times k\)，对应 \(t=1,-1\) 两个不同点。其余点满足 \(x\ne0\)，可以用 \(t=y/x\) 恢复参数。

\begin{proposition}\label{b3:prop:crossed-lines-normalization}
设 \(k\) 为代数闭域，\(A=k[x,y]/(xy)\) 为交叉直线的坐标环。限制到两轴给出嵌入
\[
A\hookrightarrow k[x]\times k[y],\qquad
\bar f\longmapsto\bigl(f(x,0),f(0,y)\bigr).
\]
其像为常数项相等的多项式对，正规化为 \(B=k[x]\times k[y]\)，导子为 \((\bar x,\bar y)_A\)，且 \(B/A\cong k\)。
\end{proposition}
\begin{proof}
\secref{b3:sec:0604}已证明限制映射的单射性与常数项相等的像，并在 \(Q(A)=k(x)\times k(y)\) 中添入 \(e=(1,0)\)，将正规化识别为 \(A[e]=B\)。原环要求两轴上的函数在交点取相同值；\(e\) 在第一条轴上取一、第二条轴上取零，正是原环所不允许的函数。正规化取消了这项值相等的限制。

若 \((f,g)\in A\) 属于导子，分别乘以 \(e\)、\(1-e\)，要求 \((f,0),(0,g)\) 的常数项仍相等，所以 \(f(0)=g(0)=0\)。反向这个条件保证与任意多项式对相乘仍在 \(A\) 内。因此导子就是两坐标常数项都为零的理想，由 \(\bar x,\bar y\) 生成。最后，差值映射 \((f,g)\mapsto f(0)-g(0)\) 的核为 \(A\)，像为 \(k\)，得到商空间的描述；\(A\) 在此商上的作用经两坐标的公共常数值给出。
\end{proof}

交叉直线的原点理想在 \(B\) 中扩张为 \(xk[x]\times yk[y]\)，因而原点的正规化纤维也为 \(k\times k\)。这里两个点分别来自两条全局不可约分量，与不可约节点曲线的两个上方点有所不同。

加厚点则用代数 \(D=k[\varepsilon]/(\varepsilon^2)\) 及其素谱描述。它只有一个素理想，却含非零的 \(\varepsilon\)。任何到域的同态都把 \(\varepsilon\) 送到零，所以所有域值点仍只能看到普通一点。若测试 \(k\) 代数自身含有非零平方零元，就可把 \(\varepsilon\) 送到这样的元，从而辨认域值点遗漏的函数变化。作为普通代数集，这一点的坐标环是 \(k\)；要研究加厚结构，必须保留 \(D\) 本身。按\cref{b3:def:normalization}的约定，本书先取 \(D_{\mathrm{red}}=k\) 再正规化，所得正规化为 \(k\)，自然映射 \(D\to k\) 的核为 \((\varepsilon)\)。这不是把 \(D\) 本身作为约化环在全分式环中取整闭包。

\subsection{约化性、不可约性与正规性}
\label{b3:subsec:0704:02}

坐标环按定义约化，所以普通代数集本身不带被点集遗漏的幂零元。不可约性进一步要求坐标环为整环。设 \(k\) 为代数闭域，\(X\subseteq k^n\) 为不可约仿射代数集。若坐标整环 \(k[X]\) 在其分式域中整闭，则称 \(X\) 为\term[zhenggui-fangshe-cu]{正规仿射簇}[normal affine variety]。正规化正是针对最后这一项条件的构造。

\begin{proposition}\label{b3:prop:normality-local-detection}
设 \(A\) 为交换含幺整环。则 \(A\) 整闭，当且仅当对 \(A\) 的每个极大理想 \(\mathfrak m\)，局部环 \(A_{\mathfrak m}\) 整闭。
\end{proposition}
\begin{proof}
正向由\cref{b3:cor:integrally-closed-localization}得到。反向设 \(z\in K=\operatorname{Frac}(A)\) 对 \(A\) 整。同一个首一方程在每个 \(A_{\mathfrak m}\) 上也成立，局部环整闭便给出 \(z\in A_{\mathfrak m}\)。\secref{b3:sec:0204}已把\cref{b3:cor:submodule-local-detection}应用于 \(A\subseteq K\)，证明共同分式域中的交式 \(A=\bigcap_{\mathfrak m\;\text{极大}}A_{\mathfrak m}\)。因此 \(z\in A\)，即 \(A\) 整闭。
\end{proof}

尖点和节点都约化且不可约，因为它们的坐标环是 \(k[t]\) 的子环。但 \(t\) 属于其分式域、满足首一方程且不在原环中，因此都不正规。交叉直线约化而可约，其正规化将两个不可约分量分开；加厚点的代数模型则首先就不约化。不能只用“有几个点”或“有几条曲线”代替这些环条件。

\ProseParagraphBreak
尖点的参数映射还给出同胚而非同构的例子。记 \(A=k[t^2,t^3]\)、\(B=k[t]\)，分别考虑谱映射 \(\pi_{\mathrm{Spec}}:\operatorname{Spec}B\to\operatorname{Spec}A\) 和经典参数映射 \(\pi(k):k\to C\)、\(t\mapsto(t^2,t^3)\)。对任一素理想 \(\mathfrak q\subseteq B\)，将\cref{b3:lem:integral-field-criterion}用于整包含 \(A/(\mathfrak q\cap A)\subseteq B/\mathfrak q\)，便知 \(\mathfrak q\) 极大当且仅当其收缩极大。在弱零点定理的识别下，\(\pi_{\mathrm{Spec}}\) 对极大谱的限制就是前面已经核验双射性的 \(\pi(k)\)。\(B=k[t]\) 的唯一非极大素理想为 \((0)\)；由整扩张谱映射的满射性，\(A\) 的非极大素理想只能是它的收缩 \((0)\)。所以两个谱的非闭点也唯一对应，\(\pi_{\mathrm{Spec}}\) 确为双射。谱映射本来连续，\cref{b3:prop:integral-spec-closed}又说明它是闭映射，因而是同胚。

经典点集上的同胚也可由闭映射直接核验。对任一理想 \(J\subseteq B\)，极大性与收缩相容，故同一闭映射公式限制为
\[
\pi_{\mathrm{Spec}}\bigl(V_B(J)\cap\operatorname{MaxSpec}B\bigr)
=V_A(J\cap A)\cap\operatorname{MaxSpec}A.
\]
因此连续的经典参数映射 \(\pi(k)\) 也是闭双射，从而为同胚；但其坐标环同态 \(A\hookrightarrow B\) 不满，因为 \(t\notin A\)，所以它不是代数集的同构。拓扑空间相同并不意味着正则函数相同。

\subsection{几何图像与局部环计算}
\label{b3:subsec:0704:03}

对尖点取 \(A=k[t^2,t^3]\)、原点极大理想 \(\mathfrak m=(t^2,t^3)_A\)。此时 \(t\notin A_{\mathfrak m}\)。否则存在 \(a,b\in A\)、\(b\notin\mathfrak m\)，使 \(tb=a\)。\(b\) 的常数项非零，所以 \(tb\) 的一次项非零；\(A\) 中元素的一次项却全为零，矛盾。这证明不整闭性确实保留在特殊点的局部环中。在导子零点集 \(V_A(\mathfrak c)\) 之外，每个素理想都避开某个导子元素，由\cref{b3:prop:conductor-basic}两侧局部化相同，因而没有这项差别。

\ProseParagraphBreak
对节点则取 \(A=k+(t^2-1)k[t]\)、原点极大理想 \(\mathfrak m=\bigl(t^2-1,t(t^2-1)\bigr)_A\)。这里同样不能把 \(t\) 写成原局部环的元素。若 \(tb=a\)、\(a,b\in A\) 且 \(b\) 在节点原点处非零，即 \(b(1)=b(-1)\ne0\)，在 \(t=1,-1\) 处分别取值得 \(a(1)=b(1)\)、\(a(-1)=-b(-1)\)。但 \(a\in A\) 要求 \(a(1)=a(-1)\)，与特征不为二矛盾。原局部环仍是整环，正规化上方却有两个点；\secref{b3:sec:0606}已具体核验其半局部正规化，并说明这两个上方点与形式方程的两条分支分别属于什么环境。

\ProseParagraphBreak
设 \(R\) 为交换含幺局部环，\(\mathfrak m\) 为其唯一极大理想，\(\kappa=R/\mathfrak m\) 为剩余域。则 \(\mathfrak m/\mathfrak m^2\) 是 \(\kappa\) 向量空间，称为该局部环的\term[yuqie-kongjian]{余切空间}[cotangent space]；其 \(\kappa\) 对偶 \(\operatorname{Hom}_{\kappa}(\mathfrak m/\mathfrak m^2,\kappa)\) 是 Zariski 切空间。在以上平面尖点、节点和交叉直线的原点，关系多项式都没有一次项，故 \(\bar x,\bar y\) 给出二维余切空间：若其线性组合落在 \(\mathfrak m^2\)，比较多项式的一次项只能使两个系数都为零。局部化后这个计算不变，因为不在原点极大理想中的元素在该商上按其非零常数值可逆作用。

这一计算只保留一阶条件。例如在 \(k[\varepsilon]/(\varepsilon^2)\) 中代入 \(x=\alpha\varepsilon\)、\(y=\beta\varepsilon\)，尖点关系 \(y^2-x^3\) 对任意 \(\alpha,\beta\in k\) 都变成零，所以 Zariski 切空间是整个 \(k^2\)。若再看方程最低次数的非零项，则得到 \(y^2\)；它所确定的切锥约化后只有方向 \(y=0\)。因此这一条方向与二维的一阶切空间并不矛盾，二者保留的阶数信息不同。

\ProseParagraphBreak
直线 \(k[t]_{(t)}\) 的余切空间却只有基 \(t+(t^2)\)，维数为一；加厚点 \(k[\varepsilon]/(\varepsilon^2)\) 的余切空间也为一维，而普通点 \(k\) 的余切空间为零。\cref{b3:tab:curve-thick-point-data}汇集了上述四个模型的数据：节点只有一个全局不可约分量，正规化在原点上方却有两个点；加厚点与尖点的集合纤维都只有一个点，纤维代数却不同。仅凭余切空间维数尚不能判断光滑性；正则局部环、Jacobian 判据及光滑性之间的联系见后面的章节。

\begin{table}[htbp]
\centering
\caption[曲线与加厚点的局部数据]{曲线与加厚点的局部数据。设 \(k\) 为代数闭域，节点例要求 \(\operatorname{char}k\ne2\)。各模型的代数记为 \(A\)，原点极大理想记为 \(\mathfrak m\)，\(A/\mathfrak m=k\)。不可约分量数指 \(\operatorname{Spec}A\) 的分量数，纤维代数为 \(\widetilde A\otimes_Ak\)，余切维数为 \(\dim_k\mathfrak m/\mathfrak m^2\)。加厚点按本书约定先约化，其正规化为 \(k\)。}
\label{b3:tab:curve-thick-point-data}
\begin{tabularx}{\textwidth}{@{}l c>{\raggedright\arraybackslash}X c c@{}}
\toprule
模型 & 不可约分量数 & 正规化纤维代数 & 纤维点数 & 余切维数\\
\midrule
尖点 & \(1\) & \(k[t]/(t^2)\) & \(1\) & \(2\)\\
节点 & \(1\) & \(k\times k\) & \(2\) & \(2\)\\
交叉直线 & \(2\) & \(k\times k\) & \(2\) & \(2\)\\
加厚点 & \(1\) & \(k\) & \(1\) & \(1\)\\
\bottomrule
\end{tabularx}
\end{table}
